% ===========================================================================
\chapter{An index}\label{ch:indexing}

Semi-naive evaluation stopped the engine repeating itself. There is a second
kind of waste left, inside every join. To satisfy \code{link(Y, Z)} with
\code{Y} already bound to \code{t7}, the join tries every \code{link} fact
and throws away all but the one that starts at \code{t7}. A relation that
could find ``the facts whose first argument is \code{t7}'' directly would
save all those attempts.

That is an \emph{index}, as a relational database keeps one: a map from the
value of the first argument to the facts that have it.

\section{An indexed relation}

Everything the join needs from a relation goes through \code{candidates},
so an indexed relation is a subclass that keeps the map up to date in
\code{add} and uses it in \code{candidates} whenever the pattern's first
argument has a value. When it does not, as in \code{route(X, Y)} with
\code{X} unbound, there is nothing to look up, and all the facts are
candidates as before.

\begin{append}{database.nex}
-- A relation that also files its facts by their first argument, so that
-- a pattern whose first argument is known looks only at the facts that
-- have it.
class Indexed_Relation inherit Relation
create
  make(p: String)
  do
    Relation.make(p)
    by_first := {}
  end
feature
  by_first: Map[String, Array[Atom]]

  add(a: Atom): Boolean
  do
    result := Relation.add(a)
    if result and a.arity > 0 then
      let key := term_text(a.args.get(0))
      if not by_first.contains_key(key) then
        by_first.put(key, [])
      end
      by_first.get(key).add(a)
    end
  end

  candidates(pattern: Atom, b: Bindings): Array[Atom]
  do
    result := facts
    if pattern.arity > 0 then
      let first := substitute_term(pattern.args.get(0), b)
      if not is_variable(first) then
        result := by_first.try_get(term_text(first), [])
      end
    end
  end
end
\end{append}

\code{Relation.add(a)} calls the version of \code{add} that \code{Relation}
declares, the one that stores the fact, before filing it in the index.

The indexed relation inherits the postcondition of \code{candidates},
\code{complete}. That is the one property an index must have, and the easy
one to get wrong: forget to file a fact under its first argument and the
join silently misses it. Here every call is checked against the full list
of facts.

\section{An indexed database}

Which kind of relation a database uses is decided in one place,
\code{new_relation}, so an indexed database is a subclass that overrides
that and nothing else:

\begin{append}{database.nex}
-- A database whose relations are indexed.
class Indexed_Database inherit Database
create
  make
  do
    Database.make
  end
feature
  new_relation(p: String): Relation
  do
    result := create Indexed_Relation.make(p)
  end
end
\end{append}

Nothing else in the engine changes. The evaluators, the join and the
queries all talk to a \code{Database}, and an \code{Indexed_Database} is
one. Extending a program by adding a subclass, rather than editing code
that already works, is sometimes called the \emph{open-closed principle}:
open for extension, closed for modification.

\section{Trying it out}

\begin{tryit}{try_indexing.nex}
intern Parser
intern Analysis
intern io/Path

function measure(label: String, db: Database, e: Evaluator)
do
  match parse_program(create Path.make(program.working_directory + "/chain.dl").read_text) of
    Ok as ok then
      do
        evaluate_program(ok.value, db, e)
        console.print_line(label + db.count.to_string + " facts, "
          + e.attempts.to_string + " facts looked at")
      end
    Err as er then console.print_line(er.error)
  end
end

measure("naive:                ", create Database.make, create Naive_Evaluator.make)
measure("semi-naive:           ", create Database.make, create Semi_Naive_Evaluator.make)
measure("semi-naive, indexed:  ", create Indexed_Database.make, create Semi_Naive_Evaluator.make)
\end{tryit}
\tryoutput{try_indexing}

The index more than halves the semi-naive evaluator's work again. Together
the two improvements look at about one fact in twenty-six of the ones
naive evaluation tries, for the same result.
